\section{Objetivos}
En esta sección se plantean los objetivos generales y específicos que guiarán el desarrollo de este trabajo de investigación.

\subsection{Objetivo general}
Desarrollar un sistema automático que, a partir de una grabación de audio de guitarra acústica, identifique el acorde que se está ejecutando con una precisión exacta mínima del 85 $\%$ y genere su representación simbólica en cifrado americano.

\subsection{Objetivos específicos}
\begin{itemize}

\item  Extraer y comparar cuatro tipos de características acústicas (transformada de Fourier de tiempo corto, CQT, chromagram y energía de armónicos) para determinar la representación más adecuada para la detección de acordes.
\item  Implementar y entrenar tres arquitecturas de aprendizaje automático (CNN 2‑D, CNN + RNN, Transformer) y optimizar sus hiperparámetros mediante búsqueda bayesiana.
\item  Evaluar el desempeño de cada modelo mediante exact‑match accuracy, top‑3 accuracy y macro‑F1, utilizando validación cruzada de 5‑fold.
\item  Analizar la robustez del mejor modelo frente a variaciones de afinación, técnicas de ejecución y presencia de ruido de fondo, reportando la degradación del rendimiento por cada factor.

\item  Documentar el proceso completo y publicar los resultados en una revista de Music Information Retrieval, liberando el código fuente bajo licencia MIT en GitHub.
\item Investigar y discutir las limitaciones del sistema propuesto, así como posibles mejoras y aplicaciones futuras en el ámbito de la transcripción musical automática y el análisis de audio.

\end{itemize}